\chapter{The preliminary arguments}
\label{chap:current-preliminaries}

This chapter follows the preliminary arguments in Section 2 of the
manuscript. In the finite group arguments, $\ell$ is prime. Brauer
characters take values in a specified field of characteristic zero,
using the chosen correspondence between roots of unity in that field
and in the field of characteristic $\ell$. Whenever a published
statement is expressed with complex characters, its application over
another value field includes the stated character and root
identifications.

\section{Central characters and tensor stabilisers}

Let $Z$ be a central subgroup of $G$ of order prime to $\ell$.
A character lies over $\nu\in\Irr(Z)$ when its restriction to $Z$
is a positive integral multiple of $\nu$. A weight $(Q,\theta)$
lies over $\nu$ when the inflation of $\theta$ to $N_G(Q)$ does.
These are the central sectors of manuscript Definition 2.1.

Over an algebraically closed field $k$ of characteristic $\ell$,
the primitive idempotents $e_\nu$ of $kZ$ are orthogonal central
idempotents of $kG$ with sum one. If $b$ is a primitive central
idempotent of $kG$, exactly one $be_\nu$ is nonzero, and then
$be_\nu=b$. This defines the sector of $b$. Every simple module
in $b$ has central character $\nu$. The ordinary characters in
the block have the same central character by their decomposition
on $\ell$-regular elements. Brauer's definition of block induction
preserves it: compare the central functions on the singleton
classes of elements of $Z$ \cite[Definition~2.1]{Spath2013}.
Thus the weight classes in $b$ also lie in this sector.

Automorphisms transport the sector of $b$ by their action on $Z$.
Inner automorphisms fix every block. Both statements are proved
on the primitive idempotents, with the same right action convention
as that used for the characters.

\begin{implementation}
The \lean{CentralSectors} namespace in
\module{ManuscriptIBAW.Preliminaries} provides the following results from ModularRep:
\lean{existsUnique_centralCharacterSector},
\lean{primitiveCentralIdempotentSector_eq_centralCharacter} and
\lean{rightMulAction_smul_centralCharacterSector}.
ModularRep constructs the multiplicative central character from the scalar
endomorphism statement of Schur's lemma in Mathlib. The central sector
deductions use these characters and the block definitions in
ModularRep together with Mathlib's group algebra and idempotent definitions.
The application to ordinary characters and weights additionally uses the specified reduction and local laws of
block induction.
\end{implementation}

For manuscript Lemma 2.2, let $N\lhd G$ and let $C$ be the group of
linear characters of $G/N$ of order prime to $\ell$, acting on
blocks by tensoring. If $\lambda\in C_b$, the common central
character of the Brauer characters of $b$ gives
$\lambda|_{Z(G)_{\ell'}}=1$. Its order makes it trivial on
$O_\ell(Z(G))$ as well. Hence it factors through $G/NZ(G)$ and
\[
 |C_b|\leq |G:NZ(G)|.
\]
The same conclusion holds for linear Brauer characters when
reduction identifies them with these ordinary characters and
commutes with tensoring. The library manual proves the result
under these coefficient and tensor compatibility hypotheses.
Its additional semidirect product argument is used later to show
that the relevant block stabilisers are $2$-hypoelementary.

\section{The full condition and modular character triples}

The full inductive condition of Sp\"ath requires a partition by
radical subgroups, equivariant local bijections, preservation of
central characters and induced blocks, and compatible extensions
over every required intermediate subgroup
\cite[Definition~4.1]{Spath2013}. For the trivial radical subgroup
it also prescribes $\chi^0\mapsto\chi$ for each ordinary character
of defect zero, with identical global and local extensions.

Manuscript Theorem 2.3 restates Sp\"ath's character triple
formulation \cite[Theorem~4.4]{Spath2017Triples} of the extension
and block compatibility requirements in
\cite[Definition~4.1(iii)]{Spath2013}. It begins on the full
universal cover $\widehat G$ of a nonabelian simple group.
Fix a radical subgroup $\widehat Q$, a Brauer character
$\widehat\theta$ of $\widehat G$, and an ordinary character of
$N_{\widehat G}(\widehat Q)/\widehat Q$ with defect zero whose
inflation $\widehat\theta'$ belongs to a block inducing to
$\bl(\widehat\theta)$. Let $\widehat Z$ be the kernel of the
central character of $\widehat\theta$ and set
\[
 G=\widehat G/\widehat Z,\qquad
 Q=\widehat Q\widehat Z/\widehat Z,\qquad M=N_G(Q).
\]
Both characters descend to these quotients. The $\ell$-part of
$\widehat Z$ lies in $\widehat Q$. Descent through its part of order
prime to $\ell$ uses the equality of the induced blocks.

These extension and block compatibility requirements are
equivalent to the required block isomorphism of the modular
character triples on $G$ and $M$. In the extension formulation,
there is a group $A$ with $G\lhd A$,
$C_A(G)=Z(A)$ and $\ell\nmid|Z(A)|$, whose conjugation image is
the required character stabiliser. The global and local Brauer
characters extend to $A$ and $N_A(Q)$, and the extensions satisfy
\[
 \bl(\widetilde\theta|_J)
 =\bl(\widetilde\theta'|_{N_J(Q)})^{\,J}
 \qquad(G\leq J\leq A).
\]
The induction superscript $J$ is essential. It is omitted in the printed equation \cite[(12)]{Spath2017Triples}.
The local triple uses the full stabiliser of the weight representative,
which must lie in the stabiliser of the global character.

\begin{implementation}
The \lean{SpathCriterion} namespace in
\module{ManuscriptIBAW.Preliminaries} provides
\lean{current_theorem_2_3} and its two implications from ModularRep.
\lean{FullCover}, \lean{PairData} and \lean{HattedPairJoins}
specify the full cover, central quotient, normaliser and
the same ordinary local character. \lean{Theorem44Source}
supplies the published equivalence uniformly for those data.
The modular block triple relation is the parameter
\lean{BlockTripleSourceSemantics.blockIsomorphic}. \lean{Theorem44Source}
assumes that, under its cover, coefficient and root hypotheses,
\lean{ClauseIII} is equivalent to \lean{TripleCondition}. The latter includes
containment of the full local stabiliser in the global stabiliser. The source
hypothesis therefore constrains the chosen relation. Its identification with
the usual definition by projective representations is not formalised, and
the transport laws used below remain assumptions. The equivalence concerns
the extension and intermediate block conditions. It does not establish the
other parts of the inductive condition. Its ordinary field is algebraically
closed, with no discrete valuation ring required.
This named theorem is a standalone formulation. The main applications
use their stated extension and block conditions and the separate
criterion sources attached to those applications. They do not invoke
\lean{current_theorem_2_3}.
In particular, \lean{SpathProposition46Theorem44Source} in
\module{ModularRep.PaperProofs.EvenFieldFLZBAWGoodFamily} assumes the passage
from a specified full family on the covering group to equivariant block
bijections satisfying the relation specified by
\lean{SpathProposition46Theorem44FamilySemantics}, under its cover and
automorphism hypotheses. This assumption constrains the relation on the
corresponding pairs. Its interpretation as block isomorphism of modular
character triples remains external. The separate
\lean{SpathProposition46Theorem44ToDefinition35FamilySource} assumes the
implication to the specified Definition~3.5 relation. These application
sources are supplied separately from the standalone theorem.
\end{implementation}

The paragraph following manuscript Theorem 2.3 applies the theorem
to a block bijection supplied by the full inductive condition.
After inflation to the full universal cover, the theorem gives the
character triple condition on the quotient by the kernel of the
central character. Inflation of this relation gives the condition
for an iBAW bijection \cite[Remark~4.4]{BroughSpath2022}.
This uses transport of a character triple relation through a central
kernel. The general transport law is an external assumption in the
formalisation. It is distinct from the central $\ell$-group argument
for manuscript Lemma 2.6 below.

\section{Combining blocks and descending to the simple group}

Manuscript Remark 2.4(a) combines Koshitani--Sp\"ath's Lemma 3.3
\cite{KoshitaniSpath2016} and Sp\"ath's Definition 5.17 and Remark 5.18
\cite{Spath2013}. The full
inductive BAW condition for one block in each automorphism orbit
gives the relative condition, and the relative conditions combine
over the required defect groups.
Taking all blocks and all radical subgroups gives the full
condition for the simple group.

Part (b) of the remark gives an elementary observation about finite
group actions. Suppose that $A$ permutes disjoint families of sets
$\mathcal U_b$ and $\mathcal V_b$ through its action on their index set.
Bijections $\Phi_b:\mathcal U_b\to\mathcal V_b$ that commute with $A_b$,
chosen for representatives of the $A$-orbits, extend to one
$A$-equivariant bijection by $\Phi(x^a)=\Phi_b(x)^a$.
This observation applies to tensor actions as well as automorphism
actions. It combines bijections and does not supply the additional
character and block conditions in the full inductive condition.

\begin{implementation}
The \lean{BlockFamilies} namespace in
\module{ManuscriptIBAW.Preliminaries} provides the ModularRep application
in which all blocks are included. Its \lean{Source} assumes the
published passage from the blockwise conditions to the full
condition, including its normalisation, over one specified
modular system. The blocks, local block
operations and root correspondence refer to that system.
When the full inductive condition holds for every block,
\lean{fullFamily} applies that hypothesis to obtain a compatible
family for all blocks. The source includes the
published passage to common choices, which can change the
individual correspondences. The reduction for the full condition in
Remark 2.4(a) is used through this special case. This provider does not
formalise the general construction in part (b). The ordinary coefficient
field is the fraction field of the specified
discrete valuation ring, with the required roots. It need not be
algebraically closed.
\end{implementation}

For manuscript Corollary 2.5, suppose that the simple group $S$
satisfies the full inductive condition at $\ell$. Sp\"ath's
Proposition 4.6, with the perfect subgroup equal to $S$, group $\Aut(S)$ and multiplicity one, supplies the partition,
block bijections and extension information
\cite[Proposition~4.6(a)--(c)]{Spath2013}.
The conversion to modular character triples uses
\cite[Proposition~3.6(b) and Theorem~3.5(b)]{Spath2017Triples}.
It gives an iBAW bijection on every block of $S$.

\begin{implementation}
The \lean{Descent} namespace in \module{ManuscriptIBAW.Preliminaries}
provides the descent construction from ModularRep. Its
\lean{FullCoverPrimeToDiagram} fixes the full cover, its quotient by the
central $\ell$-core, and the simple group. The Brauer inflation,
ordinary local characters, blocks and coefficient
comparisons refer to these maps. Lean constructs the required
epimorphism from the $\ell'$-cover to the target.
\lean{Proposition46Theorem44Source} applies the published
descent uniformly to that epimorphism and the assumed full
inductive condition. \lean{current_corollary_2_5_coherent} gives the
coherent family, and \lean{current_corollary_2_5} identifies its
triples with those in the definition of an iBAW bijection. These
interfaces use algebraically closed ordinary fields and no
modular system. Their covering maps, root comparisons and
published implications are distinct assumptions.
These named algebraically closed formulations are standalone results.
The exceptional Type B application uses the finite splitting formulation
in \module{ModularRep.CurrentSpathFiniteSplittingDescent}.
Its theorem \lean{current_corollary_2_5_block} transports the bijection
of an established covering family through the stated character and
weight correspondences. The external source supplies the character
conditions for those same pairs, under explicit splitting, root and
automorphism stabiliser hypotheses.
\end{implementation}

\section{Lifting a bijection through the central core}

For manuscript Lemma 2.6, let $G\lhd A$ and suppose that
$P=O_\ell(G)\leq Z(G)$. Write $\overline G=G/P$ and
$\overline A=A/P$. Every radical $\ell$-subgroup $Q$ of $G$
contains $P$, and we write $\overline Q=Q/P$.
Let $b$ be an $A$-stable block and let $\overline b$ be its dominated
block on $\overline G$. Suppose that $\overline b$ has an iBAW
bijection compatible with $\overline A$.

Inflation identifies the Brauer characters of these two blocks
(see, for example, \cite[Theorem~9.10]{Navarro1998}). Centrality of $P$ supplies the
hypothesis of this block correspondence. The quotient map induces
\[
 N_G(Q)/Q\cong N_{\overline G}(\overline Q)/\overline Q.
\]
Transport the weight character along this isomorphism. With the block correspondence, this gives the
weight bijection \cite[Lemma~2.1]{FengLiZhang2019}.
Both identifications commute with the action of $A$. Composing
them with the given bijection therefore gives an $A$-equivariant
bijection on $b$.

Let $(Q,\theta)$ represent the weight class corresponding to
$\chi$. The same subgroup $P$
lies in $N_G(Q)$ and in the kernels of $\chi$ and the inflated
Brauer reduction $\theta^0$. The quotient identifies the
character and weight stabilisers with their counterparts in
$\overline A$. The implication for a normal kernel in
Mart\'inez--Rizo--Rossi therefore lifts the block isomorphism of
the associated triples \cite[Lemma~3.14]{MartinezRizoRossi2026}.
This proves compatibility with $A$ for every representative of
the corresponding weight class.

\begin{implementation}
The \lean{CentralCore} namespace in
\module{ManuscriptIBAW.Preliminaries} provides the following quotient results from ModularRep:
\lean{rawLocalQuotientEquiv}, \lean{quotientRawWeight} and
\lean{ownReduction_inflation}. The last identity compares
reductions of the same ordinary local character.
\lean{lemma_2_6_compatibleBijection} constructs the bijection
on the specified block from the bijection on the quotient block,
proves equivariance, and applies \lean{sameWeight_relation_upward}.
Equivariance and injectivity also show that the stabiliser of the
weight representative is contained in the stabiliser of its Brauer
character.

\lean{CentralBlockSource} specifies the centrality hypothesis and
the block correspondence. \lean{TripleQuotientSource}
identifies the quotient stabilisers, their character values and
roots, and the block isomorphism relation in those coordinates.
The relation for the lifted characters then follows from the published
implication for a normal kernel.
\lean{returnCompatibleThroughCentralQuotient} then identifies
the constructed family with the intended blocks and
definition of an iBAW bijection. It does not require the quotient to
be a universal cover. The application over a modular system uses the stated splitting
conditions on its fraction field, without requiring algebraic closure.
\end{implementation}

\section{Permutation lattices and cyclic extensions}

For manuscript Corollary 2.7, let $A$ be $r$-hypoelementary and
let $Y,Y'$ be finite $A$-sets.
Every subgroup of $A$ is again $r$-hypoelementary. Conlon's
permutation lattice theorem \cite{Conlon1968}, in the formulation
of Bouc \cite[Theorem~3.5.5]{Bouc2000}, therefore gives $|Y^U|=|Y'^U|$ for every subgroup $U\leq A$ whenever $\Z_r[Y]\cong\Z_r[Y']$. Burnside's injectivity theorem for the mark homomorphism
\cite[Theorem~2.3.2]{Bouc2000} then gives $Y\cong Y'$ as $A$-sets.

Manuscript Lemma 2.8 applies this to an $A$-stable integral basic
set $\mathcal X$ for an $\ell$-block $b$. The decomposition map restricts to an isomorphism
\[
 d_b:\Z[\mathcal X]\longrightarrow\Z[\IBr(b)]
\]
\cite[(1.1)--(1.2)]{GeckHiss1991}.
If $d_b$ is equivariant, it is an isomorphism of permutation
lattices. Extend scalars to $\Z_2$ and apply the preceding
argument when $A$ is $2$-hypoelementary. The prime used for the
permutation lattice is $2$, independently of $\ell$. This proves
existence of an equivariant bijection. It specifies neither an
image for each character nor a unitriangular decomposition
matrix. The library manual gives the lattice and mark arguments
and the hypotheses of their formal statements.

\begin{implementation}
In \module{ModularRep.PaperProofs.ConlonBasicSet},
\lean{corollary_2_4_conlonMark} implements manuscript Corollary 2.7.
It deduces an equivariant bijection from the permutation lattice
isomorphism, assuming \lean{PadicConlonMarkDetection} and
\lean{PublishedBurnsideMarkInjectivity}. These hypotheses state the two published results for arbitrary finite $A$-sets.
The same module's \lean{lemma_2_5_basicSetBridge} implements the
lattice argument in manuscript Lemma 2.8. Starting from an integral
linear equivalence for the restricted decomposition map and its
equivariance, it constructs the permutation lattice isomorphism,
extends scalars to $\Z_2$, and applies the preceding deduction.
The existence and equivariance of the restricted decomposition map
remain hypotheses of this declaration.
\end{implementation}

Manuscript Lemma 2.9 is the cyclic extension theorem. If $N\lhd H$
and $H/N$ is cyclic, an invariant irreducible ordinary or Brauer
character of $N$ extends to $H$. For examples of these results, see \cite[Corollary~11.22]{Isaacs1976} for ordinary characters and \cite[Theorem~8.12]{Navarro1998} for Brauer characters.
The ordinary statement uses complex characters
or a compatible splitting interpretation. The modular statement
uses algebraic closure of the field of characteristic $\ell$
and requires no coprimality condition on $H/N$.

\begin{implementation}
\module{ModularRep.OrdinaryCharacterCyclicExtensionBridge} proves
\lean{exists_extension_of_fixed_cyclic_quotient}, using
\lean{Representation.CyclicExtensionPrinciple}.
\module{ModularRep.BrauerCharacterExtensionBridge} proves
\lean{exists_extension_realisation_of_ibr_fixed_cyclic_quotient},
using \lean{Representation.BrauerCyclicExtensionPrinciple}.
The Brauer character is then constructed in
\module{ModularRep.BrauerCharacterHomPullback} by
\lean{Representation.Extension.brauerCharacterExtensionWitnessOfCompatible},
with compatible root correspondences.
These constructions prove invariance of the chosen realisations and the
character restriction identities. The representation extension principles
remain explicit hypotheses. The ordinary principle over the chosen field
quantifies over all finite ambient groups. Splitting for the base group alone
does not establish it. Isaacs's theorem gives the complex case, and its
interpretation over the chosen ordinary field is an additional assumption.
The modular principle uses an algebraically closed field. Neither extension
principle is proved in these declarations.
\end{implementation}

Finally, manuscript Lemma 2.10 assumes that the simple group $G$
is its own universal cover and that $\Aut(G)=G\rtimes E$
with $E$ cyclic. For an $E$-stable block $b$, an $E$-equivariant
bijection $\IBr(b)\to\Alp(b)$ extends to equivariance under
$G\rtimes E$, since inner automorphisms fix both sets. Each
Brauer character extends to its stabiliser by the cyclic
extension theorem. For a weight $(Q,\theta)$, set
$D=(G\rtimes E)_{Q,\theta}$. The quotient
$D/N_G(Q)$ embeds in $E$. Apply the same extension theorem to
the reduction of the original local character on $N_G(Q)/Q$
inside $D/Q$, then inflate to $D$.

These extensions and the block equality attached to the weight
verify the specialised Brough--Sp\"ath criterion with
$\widetilde G=G$ \cite[Theorem~3.18]{FengLiZhang2022Jordan}.
The conclusion is that $b$ is BAW-good. Applications use the
specified cover, local character, block induction and extension
interpretations. A full family conclusion also uses the separate
passage from the block conditions to compatible correspondences.

\begin{implementation}
In \module{ModularRep.PaperProofs.CurrentCyclicOuterBAW},
\lean{lemma_2_10_bawGood} applies the published criterion, supplied
as \lean{FLZ318BAWGoodSource}, to the hypotheses constructed by
\lean{completedExplicitPackage_of_cyclicEndgame}.
Its conclusion is BAW-goodness of the specified block.
The extension principles, compatible roots, structural and block
assumptions, and identifications of the covering group and character
triple relation remain explicit. This declaration does not prove the
published criterion or the separate passage to a compatible family
for all blocks.
\end{implementation}

\section{Character selection for finite reductive groups}

The manuscript proves two character selection lemmas for connected
reductive groups. Lemma 2.11 concerns isolated families in Type B.
Lemma 2.12 concerns selected Type C families in rank at least three.
Their proofs are in Appendix B of the manuscript.
For the latter, the semisimple parameter is chosen in a split torus
and multiplied by a central element to obtain a parameter whose square is one.
The family in the conformal symplectic group is attached to this
rational pair. Both lemmas select ordinary characters whose positive
Alvis--Curtis duals have Kronecker delta scalar products with the
generalised Gelfand--Graev characters of the rational classes in the
chosen geometric unipotent class. The arguments and their source
assumptions are discussed in Chapters~\ref{chap:current-type-b}
and~\ref{chap:current-type-c}.

\enlargethispage{2\baselineskip}
\begin{implementation}
The geometric constructions, including the central choice in Lemma 2.12,
are not formalised. The character theoretic identities used in the finite
proofs are explicit assumptions. The Type B application proves the finite spanning,
independence, coverage and selection deductions from those identities.
In \module{ManuscriptIBAW.TypeC.GGGRSelection},
\lean{delta_selection} proves a conditional form of Lemma 2.12 on the
specified conformal symplectic group. From the stated component action,
weighted formula, denominator equality and trace identities, the proof
shows that the component indices are one and the scalar product row sums
equal one. It then constructs an injective choice of characters. The
selected family and its geometric interpretation are supplied separately.
The general finite arguments are described in the library manual's
chapter on character counts from finite matrices.
\end{implementation}
